\documentclass{article}
\title{名スピーチ講義 — 語りの技術を読む}
\author{TeX 図書館}

\begin{document}

\section{演説とは何か}

演説は、読むための文章ではない。声に出され、時間の中で流れて消える言葉である。読み返すことができない。だから演説の技術は、書き言葉の技術とは別のものになる。この本は、パブリックドメインの名演説を材料に、その技術を分解する講義である。

\subsection{声・時間・聴衆の3条件}

演説を演説たらしめる条件は3つに集約される。声、時間、聴衆である。

\begin{itemize}
\item \textbf{声}: 音声には抑揚と間がある。同じ文でも、どこを強く読むかで意味が変わる。
\item \textbf{時間}: 言葉は一方向に流れる。聞き飛ばした一文は、二度と帰ってこない。
\item \textbf{聴衆}: 複数人が同時に聞いている。誰か一人の理解ペースに合わせられない。
\end{itemize}

この3条件が、すべての技術の土台になる。書き言葉で許される「後で戻って読めばいい」が、音声にはない。だからこそ、順番と短さが命になる。

\begin{quote}
\textbf{【定義】}\ \textbf{演説}とは、特定の聴衆に向けて、声と時間の中で行われる一方向の語りである。文字に起こされたものは記録であり、演説そのものではない。記録を読んで感動しても、その感動の半分は「場」と「声」の仕業である。
\end{quote}

\subsection{書き言葉と話し言葉の分岐}

書き言葉は読み返せる。だから修飾を重ねてよい。主語を後ろに置いても、読者は戻れる。話し言葉は読み返せない。だから一回で届く文だけが届く。息が続く長さの文しか持たない。

\begin{quote}
\textbf{【例】}\ 悪い例: 「本プロジェクトにおいて想定される最大のリスクは、人的リソースの不足であると考えられる。」 / 良い例: 「一番の不安は人手です。今の人数では、この計画は回りません。」
\end{quote}

内容は同じである。悪い例は読者を待たせる。良い例は一息で着地する。演説では、待たされた聴衆は席で考え始める。考え始めた聴衆は、次の文を聞かない。

もう1つ、分岐が大きい。主語の扱いである。書き言葉では主語を落とすと、誰がやったのかが消える。話し言葉では、主語を落としてよい。その場にいる誰が話しているかは、聴衆が知っているからである。逆に、話し言葉で丁寧に主語を補い続けると、地の文がくどくなる。

\begin{quote}
\textbf{【例】}\ 悪い例: 「私は、本日ご参加いただいた皆様方に向けて、簡単にですがご説明をさせていただきます。」 / 良い例: 「5分で、3つの話をします。」
\end{quote}

\subsection{この本の読み方}

この本は9章からなる。冒頭、構成、反復、具体、比喩、間、そして危うい技術と応用までを、名演説の断片を切り口にして進める。引用は鍵となる短い句だけを行い、長い部分はすべて日本語で言い直す。原文の暗唱が目的ではない。技法の分解が目的である。

材料はすべてパブリックドメインである。1775年のヴァージニアの議会、1933年の就任式、1940年の下院、1961年の就任式、1951年の議会。二百年前の青年と、冷戦の初年度の首相級が、同じ道具を使っている。時代も国も違うのに使える技術は、道具であるからこそ時代を超える。

\begin{quote}
\textbf{【例】}\ 同じ内容の2つの使い方。悪い例: 原稿を黙って読み上げる。/ 良い例: 原稿の3行目で一度上げて、聴衆を見てから続ける。前者は記録の再生であり、後者が演説である。
\end{quote}

\begin{quote}
\textbf{【注意】}\ 章によっては、声に出して読むと分かる箇所がある。第4章と第7章は、目で追うだけでは半分しか分からない。手を止めて、一文だけ音読してから次へ進むほうが、速く身につく。
\end{quote}

\section{冒頭の一文}

聴衆は3秒で判断する。この話は聞く価値があるかどうかを、頭の中の一秒で決めてしまう。冒頭の一文の仕事は、その判断を「聞こう」に反転させることである。

\subsection{型は3つ}

冒頭の型は、大きく3つに分かれる。

\begin{itemize}
\item \textbf{問いかけ}: 聴衆に答えてもらう。答えさせた瞬間から、聴衆は当事者になる。
\item \textbf{断言}: 事実や結論を先に置く。意外であるほど、引き込まれる。
\item \textbf{宣言}: 自分の立場と覚悟を述べる。この先に何を望むかが、最初に見える。
\end{itemize}

\subsection{ヘンリーの冒頭：褒めてから踏み込む}

パトリック・ヘンリーは1775年、ヴァージニア植民地の議会で演説した。世に知られる「Give me liberty, or give me death!」は、実は演説の結びにある。冒頭は逆だった。直前に反対意見を述べた者たちを、愛国心も能力も高いと持ち上げ、その上で自分の意見は必ず反対されると予告した。

褒めてから反論する。これは、聴衆の防御を下ろしてから本題に踏み込む型である。聴衆が自分たちを褒められた直後に攻撃されると、防御が間に合わない。ヘンリーは、技術的に文句の出ない礼を述べてから、段階的に温度を上げていった。

\begin{quote}
\textbf{【例】}\ 悪い例: 「えー、本日はお話の機会をいただきありがとうございます。私はあまり上手ではありませんが、頑張ります。」 / 良い例: 「会議は毎週90分、無駄になっています。今日の30分で、その半分を取り戻します。」
\end{quote}

悪い例の仕事は、自分の地位の低さを聴衆に伝えることだけである。良い例は、聴衆が得をする予告を1文で渡している。

\subsection{ケネディの冒頭：礼のあとに一度、着地する}

1961年のケネディ就任演説は、儀礼の名乗りから始まる。議長、副大統領、閣僚、議員と役職を並べたあとで、建国の世代の話を一度置いて、初めて本題に入る。有名な「Ask not what your country can do for you」も、その先にある。

構造は3段である。礼、共有、依頼。全員がすでに知っている記憶(独立という起源)を一度手渡してから、まだ誰もしていない依頼を出す。いきなり依頼から始めると、聴衆は関係を確認する時間を使う。その時間を、ケネディは名乗りと共有で済ませている。

\begin{quote}
\textbf{【注意】}\ 冒頭で自己紹介と前置きに全時間を割かない。「あまり上手ではないので」「念のため」と、話し手の保身を先に置くと、聴衆は安心して聞くのをやめる。保身はメモに書く。冒頭は、聴衆の得に使う。
\end{quote}

3つの型は、同時に使える。問いかけで引き、断言で自分の立場を示し、宣言で約束する。15秒の冒頭なら、この順で十分である。

\begin{quote}
\textbf{【演習】}\ 社内勉強会の冒頭一文を、問いかけ・断言・宣言の3つの型でそれぞれ1つずつ書き分けよ。テーマは「議事録の書き直し」とする。
\end{quote}

\textbf{解答の指針:} 問いかけなら「先週の決定事項を、全員が言えるでしょうか」。断言なら「議事録が書かれていない議論は、来週もう一度やることになります」。宣言なら「今日の30分で、議事録の様式を1つに決めます」。3つは同じ情報でも、聴衆の置き方が違う。問いかけは聴衆を主語に、断言は事実を主語に、宣言は話し手の約束を主語にする。

\section{構成：結論をどこに置くか}

話し言葉の構成は、書き言葉より単純でなければならない。聴衆は目次を見られないからである。使える型は、実質3つある。結論の先出し、三段構成、そして就任演説の定石である。

\subsection{結論は先か、最後か}

書き言葉では逆三角形が基本だった。結論を先に置き、根拠を重ねる。演説にも同じことが言えるが、少し違う使い分けがある。

\begin{itemize}
\item \textbf{結論を先に置く}: 時間のない場、すでに敵意のある場。待たされると離脱する聴衆に有効である。
\item \textbf{結論を最後に置く}: 情報を新しくする場、感情を共有する場。聴衆と一緒に温度を上げてから、着地点を渡す。
\end{itemize}

\begin{quote}
\textbf{【注意】}\ どちらを選ぶかは、聴衆があなたを待っているかどうかで決まる。待ってくれるなら最後、待ってくれないなら先。曖昧なまま中途半端に置くのが最悪である。
\end{quote}

\subsection{三段構成：現状・主張・要請}

多くの説得の演説は、3段で動く。

\begin{enumerate}
\item \textbf{現状}: いま何が起きているか。聴衆がすでに感じていることを、言葉にする。
\item \textbf{主張}: だからどうすべきか。事実への解釈を1つ渡す。
\item \textbf{要請}: 誰が、何をするか。行動がなければ、演説は感想で終わる。
\end{enumerate}

現状を飛ばして主張すると、同意の土台がない。要請を落とすと、感想になる。3つが揃ったとき、聴衆は帰る場所を持つ。

\begin{quote}
\textbf{【例】}\ 悪い例: 「セキュリティは重要です。ぜひ協力をお願いします。」 / 良い例: 「先月、社外への誤送信が3件ありました。現状では、送信前の確認は一人の直感です。明日から、送信前に添付と宛先を必ず2人で確認してください。」現状に数字、主張に解釈、要請に動詞が入っている。
\end{quote}

三段の診断は簡単である。聞き返されたときに「え、何をすれば?」と聞かれたら、要請が抜けている。「で、結局どうしたいの?」と聞かれたら、主張が抜けている。「そんなの知ってるよ」と返されたら、現状が自分本位である。

\subsection{就任演説の定石：危機と希望の同居}

就任演説には決まりごとがある。就任という事実の確認、その国の抱える危機の提示、そして国民への依頼。この3点は、ほぼどの就任演説にも現れる。

1933年のルーズベルト第1次就任演説は、大恐慌の只中で行われた。銀行が相次いで閉まっていた時代である。そこで彼は「the only thing we have to fear is fear itself」と断言した。国民が本当に戦うべき対象を、外部の困難ではなく「恐れそのもの」に置き換えた。次の句で、その恐れを「名前もなく、非合理的で、正当化のない恐怖」と分解している。抽象語の「恐れ」を、具体的な状態に落としたのである。

ケネディの「Ask not what your country can do for you」は、依頼の反転である。それまでの「国が何をしてくれるか」という問いを一度裏返し、国民を主語に置き直した。国民を当事者にすることが、就任演説の要請の典型形だと言ってよい。

\section{反復とリズム}

声に出したときに効く技法は、紙で読むと単調に見える。反復がそうだ。単調さは、音声では安心になる。同じ形の文が来ると、聴衆は次を予測できる。予測できた瞬間、聴衆は安心する。

\subsection{アナフォラ：文の頭を揃える}

アナフォラは、複数の文の頭に同じ語を置く技法である。チャーチルは1940年6月、下院で連続して語った。「we shall fight」という頭を何度も反復し、海、上陸地点、丘陵、街道と場所をずらしていく。頭が同じだから、聴衆は最後の一行まで足場を失わない。変化するのは場所だけで、姿勢は変わらない。この落差が、決意の厚みとして届く。

\subsection{三項並列：3つで一段落を閉じる}

三項並列は、事項を3つ並べて締める形である。3つは多いすぎず、少なすぎない。2つだと対比になり、4つ以上だと列挙になってしまう。ルーズベルトの「名前もなく、非合理的で、正当化のない恐怖」は、形容詞の三項並列である。恐怖という抽象語に、3つの形容詞で形を与えた。

\begin{quote}
\textbf{【例】}\ 悪い例: 「この施策は、効率を高め、コストを下げ、品質を維持し、社員の満足度を向上させ、顧客対応を改善する。」 / 良い例: 「この施策は、速く、安く、壊さない。この3つを同時に狙います。」
\end{quote}

悪い例は5項並列である。どの項目も残らない。良い例は3項で、しかも価値の衝突(速さと価格と品質)を1文に畳んでいる。

\begin{quote}
\textbf{【注意】}\ 反復は、中身が変わらなければ効かない。頭を揃えて中身を並べるだけでは、単調が単調を生む。チャーチルの反復では、場所だけが変わっていた。変化する部分を、必ず1つ入れる。
\end{quote}

\subsection{一文は息の長さで切る}

話し言葉の一文は、息が続く長さで切る。目安は、書き言葉の一文より短いほうである。同じ語尾が3つ続くと、聴衆は音だけを聞くようになる。文を変えるたびに、意味が戻ってくる。

\begin{quote}
\textbf{【演習】}\ 「先月の問い合わせ対応の平均応答時間は前月比で大きく改善したが、一方で一次対応の担当者数が減っているという課題も残っており、両方を見た改善策を提案したい」という一文を、声に出して短い文に割れ。
\end{quote}

\textbf{解答の指針:} 判断を3つに割る。「応答時間は前月より短くなりました。」「しかし一次対応の担当者は減っています。」「両方を見る改善案を出します。」割ったあと、音読して息切れする箇所がなければ合格である。数字は、割っても残る。削るのは文の長さであり、情報ではない。

\section{具体と数字}

抽象語は、聴衆に翻訳を任せる言葉である。翻訳を任せられた聴衆は、勝手にイメージを作る。作ったイメージが話し手の意図と一致する保証はない。具体語は翻訳済みである。演説では、この差が信頼の差になる。

\subsection{抽象語は何かで支える}

「努力」「成長」「安定」「危機」。こうした語は、それ単独では聴衆に何も渡さない。裏を何かで支える必要がある。支える材料は3つある。数字、場面、身体語である。

\begin{itemize}
\item \textbf{数字}: 「多くの問い合わせ」は「月400件」に置き換える。
\item \textbf{場面}: 誰が、どこで、どう困っているかを1つ描く。
\item \textbf{身体語}: 「苦労」は「血と汗と涙」に置き換える。
\end{itemize}

\subsection{チャーチルの具体：抽象語を身体に置き換える}

1940年5月、チャーチルは首相として最初の演説でこう述べた。「I have nothing to offer but blood, toil, tears and sweat」。直訳は、血、労苦、涙、汗、の4つである。

注目すべきは、苦労という抽象語を、血や涙や汗という身体語に置き換えた点にある。苦。苦労は聴衆の一人一人で意味が違う。血と汗は、誰でも同じ絵を描く。さらに彼は、自分の提供物を「持ち合わせはこれしかない」と限定して述べた。誇張ではなく、限定の形をとった約束は、逆に信頼される。

\begin{quote}
\textbf{【例】}\ 悪い例: 「我々は今、非常に厳しい局面において、大きな努力を要求されている。」 / 良い例: 「必要なのは、血と汗と涙です。」
\end{quote}

同じ内容を支える材料が違うだけで、聴衆が思い浮かべる絵はまるで違う。前者は会議室、後者は戦場である。

\subsection{ルーズベルトの具体：恐れを分解する}

ルーズベルトは「恐れそのもの」と断言したまま、そこで終わらせなかった。「名前もなく、非合理的で、正当化のない恐怖」と、恐怖に3つの形容詞を付けた。恐怖という1語を、実在しない名前、役に立たない思考、根拠のない怯え、に分解した。分解した瞬間、恐怖は「扱えないもの」から「扱えるもの」に変わる。

\begin{quote}
\textbf{【注意】}\ 数字は正確であること。曖昧な大げさな数字より、控えめで正しい数字のほうが聴衆を動かす。出所を言えない数字は、聞かれたら壊れる。壊れる数字は、入れないほうがよい。
\end{quote}

\begin{quote}
\textbf{【演習】}\ 「サービスの利用が伸びており、社内でも導入の価値が認められるようになってきている」という一文を、数字・場面・身体語のいずれかで支えて書き直せ。
\end{quote}

\textbf{解答の指針:} 数字なら「利用は月30件から240件になりました」。場面なら「火曜の朝、経理の3人が順番にこの画面を開いていました」。身体語なら「手作業の締め作業が、毎月朝まで粘っていたのを、今は昼前に終わります」。どれを選ぶかは聴衆次第だが、選ばない限り、元の一文は感想のままである。

\section{比喩と物語}

比喩は、聴衆がすでに知っている領域に、まだ知らないものを運ぶ技術である。物語は、抽象を一回きりの場面に変える技術である。どちらも、説明の代わりにはならない。説明を速くするものである。

\subsection{譬喩の設計}

良い譬喩は3条件を満たす。聴衆が知っていること、話し手の主張との構造が似ていること、1段落に1つであることである。条件を外すと、混乱を生む。

\begin{quote}
\textbf{【例】}\ 悪い例: 「この新機能は、ピラミッドのような、かつてないアーキテクチャです。」 / 良い例: 「この新機能は、操作を3ステップに畳みました。コーヒーを淹れる時間より短いです。」
\end{quote}

ピラミッドは、聴衆の職種によって意味が変わる。コーヒーの時間は、ほぼ全員が共有している。譬喩は、話し手が気に入ったかどうかでは選ばない。聴衆が持っているかで選ぶ。

\subsection{引用と身振り}

引用は、聴衆がすでに知っている言葉に、自分の論を預ける技術である。マッカーサーは1951年の退役演説の結びで、古い歌の一節「Old soldiers never die」を引用して、それに「they just fade away」を付け足した。古い歌を知っている聴衆は、歌の感情をそのまま持ち込む。話し手は、感情を作らずに借りられる。

ガンディは「My life is my message」と短く言い切った。行動そのものをメッセージに置き換えている。身振りとは、比喩を言葉の外で実演することである。何を語るかと同じだけ、何をするかも語る。

\subsection{聴衆を登場させる}

「あなたたち」「我々」。この2つの語で、聴衆は文の主語になる。聴衆が主語に入った瞬間、演説は他人の話から自分の話に変わる。ただし、「我々」を安易に使うと、話し手と意見が違う人まで巻き込む。巻き込みが外れた「我々」は、逆効果である。

\begin{quote}
\textbf{【注意】}\ 譬喩は、当てはまらなかった瞬間に信頼を失う。聴衆の職業、年齢、経験が分からないまま、話し手だけが知っている例えを持っていくのは危険である。迷ったら、毎日の動作(淹れる、送る、並べる)から選ぶ。派手な例えより、当たり前の例えのほうが届く。
\end{quote}

\begin{quote}
\textbf{【演習】}\ 「社内の情報共有が不十分である」という一文を、(1)譬喩で、(2)1つの場面で、書き直せ。
\end{quote}

\textbf{解答の指針:} (1)は「必要な資料が、部署ごとの金庫に分けて入っている」のように、読者が知っている動作に置き換える。(2)は「月曜の朝、AさんがBさんに3回同じ質問をしている」のように、誰がどこで困っているかを1つ描く。抽象1つが、絵2つに変われば合格である。

\section{間と声}

音声には、文字にない語彙が2つある。間と抑揚である。これは朗読の技術であると同時に、構成の技術でもある。どこで間を取るかは、どこで意味の区切りを作るかだからである。

\subsection{沈黙が作る意味}

間には3つの役割がある。区切り、強調、そして待つことである。区切りの間は、段落の切れ目として働く。強調の間は、その直後に来る言葉の重さを予告する。待つ間は、聴衆に考える時間を渡す。考えさせたいときに話し続けるのは、逆の行為である。

\begin{quote}
\textbf{【例】}\ 悪い例: 「つまりこのことはつまり非常に重要なことでしてつまり皆さんの理解が大事です。」 / 良い例: 「これは、3か月で終わりません。(間) 皆さんの判断が要ります。」
\end{quote}

悪い例は、詰めることで熱量を作ろうとしている。熱量は詰め物では作れない。良い例の間は、次の一句のための助走である。1秒の沈黙は、10の形容詞に勝る。

\subsection{区切りは書きながら決める}

間は、暗唱してから足すものではない。下書きの段階で、読点と段落の切れ目に先に刻んでおく。区切りを先に決めると、速さの設計は半分済んでいる。区切りを決めないまま速さだけ整えると、意味の切れ目と声の切れ目がずれ、聴衆は息継ぎの場所で言葉を落とす。

\begin{quote}
\textbf{【例】}\ 悪い例: 「明日から朝会を30分にします理由は前日の持ち越しを減らすためです」 / 良い例: 「明日から、朝会は30分です。(間) 持ち越しを減らすためです。」
\end{quote}

\subsection{速さと音量の設計}

速さは緊張、音量は決意、に対応する。詳しい事情は速く、重要な判断は遅く。軽い背景は小さく、約束は大きく。この対応を決めてから稽古に入ると、暗記は半分で済む。順番に暗記すると、速さの設計が抜け落ちる。

\begin{quote}
\textbf{【注意】}\ 音読は推敲の代わりになる。つっかえる箇所は、そのまま聴衆もつっかえる。息が尽きる箇所は、文が長すぎる証拠である。紙で上手に見える文が、声では届かないことは珍しくない。
\end{quote}

\section{危うい技術：煽りとスローガン}

ここまでの技術は、中立である。反復も、比喩も、間も、正しい内容にも嘘にも使える。最後に扱うのは、それを見抜く力である。使う技術と、防ぐ技術は、同じ解剖学の両面である。

\subsection{煽りは反復の悪用}

煽りの典型的な形は3つある。敵の固定、規模の誇張、そして反復である。誰が悪いのかを一度決め、数を盛り、短い句で繰り返す。反復が安心を生む技術だったのは、予測できることが安心だったからだ。予測できない内容を反復すると、安心ではなく刷り込みになる。

\subsection{対立の作り方}

対立は、分類を2つに絞ることで作られる。「我々か、彼らか」。この分類に、中間も例外も入れない。名演説が対立を扱うとき、たいてい敵を人ではなく状態に置く。ルーズベルトの敵は他国民ではなく「恐れ」だった。敵が状態であるなら、聴衆は敵と和解できる。敵が人であるなら、和解の道が消える。

\subsection{スローガンの罠}

短い句は、反復できる代わりに、中身が薄い。逆に言えば、中身を確認できない句ほど、反復に耐える。「日本を取戻す」「もう一度、偉大に」。何を取り戻すのか、誰にとって偉大なのか。答えが空のまま反復された句は、聞き手の願いを勝手に吸い寄せる。

\begin{quote}
\textbf{【注意】}\ 修辞を見抜くチェックは3つである。(1) 主語は誰か。(2) 数字の出所はどこか。(3) 敵は人か、状態か。この3つに答えられなかったら、感動の前に一度、止まる。
\end{quote}

\begin{quote}
\textbf{【例】}\ 煽りの文を直す。悪い例: 「また別の誰かが、私たちの職を奪おうとしている。だから、今すぐ立ち上がるべきだ。」 / 良い例: 「来年、この職種の求人は12件減りました。原因は3つあります。そのうち1つだけ、今日ここで直せます。」前者は主語が曖昧で、数がなく、敵が人である。後者は主語を数字に移し、敵を原因に落としている。
\end{quote}

反復が善用される例もある。1963年のワシントン行進で語られた演説は、同じ一つの言葉を何度も返しながら、まだ実現していない将来を一段ずつ具体的にしていった。同じ技法が、積み上げにも刷り込みにもなる。分けるのは、句の中身が確かかどうかである。

\section{自分の一分に移す}

名演説の技術は、講壇の上にだけ存在しない。メールの1行目、会議の30秒、入場時のあいさつ。短い場面に、同じ骨格はそのまま入る。この章は総合である。

\subsection{メール：冒頭の一文がすべて}

メールの冒頭一文は、演説の冒頭である。3秒で判断されるのは同じだ。件名と1行目で、用件と依頼が見えなければ、本文は読まれない。

\begin{quote}
\textbf{【例】}\ 悪い例: 「先日ご依頼いただいた件につきまして、一部確認事項が生じましたのでご連絡いたしました。」 / 良い例: 「契約書の納期だけ、ご確認ください。金曜17時までに返事がなければ、そのまま送付します。」
\end{quote}

\subsection{会議：三段構成の30秒}

会議での発言は、現状・主張・要請の3段で組む。「今こういう状況です。だからこうしたい。だから○○さん、△△をお願いします」。30秒で終わる。三段のどれかが欠けると、議論は場当たりになる。要請が欠けると、決定事項が生まれない。

\subsection{ウェルカムスピーチ：間を先に設計する}

入場時のあいさつは、拍手と間の設計が8割である。拍手が収まるまで待つ。待てば、聴衆は席に落ち着く。落ち着いたあとに最初の一文を置くと、その一文は全員に届く。拍手の最中に話し始めると、自分の声を自分の足音に踏み潰される。

どんな短い場面でも、準備の順番は同じである。まず冒頭の一文、次に結論の置き場所、最後に声の設計。暗記はこの後である。順番を守らないと、うまく暗記した原稿を、正しくない順で話してしまう。

\begin{quote}
\textbf{【例】}\ 悪い例: 「本日はご多忙の中お集まりいただき、誠にありがとうございます。私めといたしましては、僭越ながら、本日お話をさせていただきます。」 / 良い例: 「今日ここに来られた方は、だいたい1つの疑問を持っています。この場所で、何が1つ分かりますか。答えは3つです。」
\end{quote}

\subsection{送る前のチェックリスト}

自分の発話や文を出す前に、上から順に確かめる。全8項目、所要は1分である。

\begin{itemize}
\item 冒頭の一文は、相手の得を予告しているか
\item 結論は、先か最後のどちらかに振り切っているか
\item 現状・主張・要請の3段が、揃っているか
\item 同じ語の頭を揃えた反復が、1つは入っているか
\item 三項並列で締める箇所があるか。4項以上に膨らんでいないか
\item 抽象語に、数字・場面・身体語が伴っているか
\item 譬喩は、相手が知っているものか。1段落に1つではないか
\item 短い句を繰り返しているとき、その中身は確かか
\end{itemize}

\begin{quote}
\textbf{【総合演習】}
\begin{enumerate}
\item 次の部署移動のお知らせ一文を、三段構成(現状・主張・要請)に組み直せ。「部署移動に伴い、10月より席が変わりますのでご了承ください。」
\item 30秒の会議発言を、冒頭の一文から書き下ろせ。テーマは「定例会議を60分から30分にすること」。
\item 自分の最近の発言1つを、チェックリスト8項目で自己採点せよ。
\end{enumerate}
\end{quote}

\textbf{解答の指針:} 1. 現状「AチームとBチームが別々の会議を続けています」、主張「10月から両方を1つの定例にまとめます」、要請「持ち越し議題は金曜までにチケットへ」。2. 冒頭は「この会議は毎週、30分余っています」のように数字の断言で始める。要請は「今日、参加者を6人に絞る判断を先にお願いします」まで落とす。3. 埋まっていた項目を数え、欠けた項目を次回は下書きの段階で先に決める。短い場面ほど、型の欠落はそのまま伝わりの欠落になる。

\end{document}
